% Part: computability
% Chapter: computability-theory
% Section: halting-problem

\documentclass[../../../include/open-logic-section]{subfiles}

\begin{document}

\olfileid{cmp}{thy}{hlt}
\olsection{停机问题}

根据构造，通用部分可计算函数~$\fn{Un}(e,x)$ 有定义，当且仅当由~$e$ 编码的函数对输入~$x$ 的计算产生值。自然会产生一个问题：我们能否判定这种情况是否发生？事实上，不能。对于Turing机计算模型而言，这意味着，一台给定的Turing机在给定输入上是否停机，是计算上不可判定的。因此，下述定理被称为“停机问题的不可判定性”。下面将给出两个证明。第一个延续先前的讨论，第二个则更为直接。

\begin{thm}
\ollabel{thm:halting-problem}
设
\[
h(e, x)  =
\begin{cases}
1 & \text{若\/ $\fn{Un}(e, x)$ 有定义} \\
0 & \text{否则}
\end{cases}
\]
则 $h$ 不可计算。
\end{thm}

\begin{proof}
假设 $h$ 可计算。我们将证明，由此便可定义一个通用可计算函数。定义
\[
\fn{Un'}(e,x) =
\begin{cases}
\fn{Un}(e,x) & \text{若 $h(e,x) = 1$} \\
0 & \text{否则}
\end{cases}
\]
此时 $\fn{Un'}(e, x)$ 是全函数，且若~$h$ 可计算，则它也可计算。例如，我们可以用原始递归定义 $g$：
\begin{align*}
g(0, e, x) & \simeq 0 \\
g(y+1, e, x) & \simeq \fn{Un}(e,x);
\end{align*}
进而
\[
\fn{Un'}(e,x) \simeq g(h(e,x),e,x).
\]
由于在 $\fn{Un}(e,x)$ 有定义之处，$\fn{Un'}(e,x)$ 都与其一致，因此 $\fn{Un'}$ 对恰好为全函数的那些部分可计算函数是通用的。但这与 \olref[nou]{thm:no-univ} 矛盾。
\end{proof}

\begin{proof}
假设 $h(e,x)$ 可计算。定义函数 $g$ 如下：
\[
g(x) =
\begin{cases}
  0                & \text{若 $h(x,x) = 0$} \\
  \fundefined & \text{否则}
\end{cases}
\]
函数 $g$ 是部分可计算的。例如，可将其定义为 $\umin{y}{h(x,x) = 0}$。因此，存在某个~$e$，使得对每个~$x$ 都有 $g(x) \simeq \fn{Un}(e,
x)$。$g$ 在 $e$ 处是否有定义？若有，则根据~$g$ 的定义，$h(e,e) = 0$（$h$~只有在有定义时才可能取值~$0$）。而由~$h$ 的定义，这意味着 $\fn{Un}(e, e)$ 无定义。根据我们对每个~$x$ 都有 $g(x) \simeq \fn{Un}(e, x)$ 的假设，可知 $g(e)$ 无定义，矛盾。
另一方面，若 $g(e)$ 无定义，则 $h(e,e) \neq 0$，从而 $h(e,e) = 1$。由此可得 $\fn{Un}(e, e)$ 有定义。但由于 $g(x) \simeq \fn{Un}(e, x)$，那么 $g(e)$ 亦应有定义。再次得到矛盾。
\end{proof}

\begin{tagblock}{TMs}
\begin{explain}
我们可以用Turing机的语言来描述这个论证。假设存在一台Turing机~$H$，它以Turing机~$E$ 的描述和输入~$x$ 作为输入，并判定 $E$ 在输入~$x$ 上是否停机。那么我们就能构造另一台Turing机~$G$，它接受单一输入~$x$，运行 $H$ 以判定指标为~$x$ 的机器 $M_x$ 在输入~$x$ 上是否停机，然后做出相反的动作。换言之，若 $H$ 报告 $M_x$ 在输入~$x$ 上停机，$G$ 就进入无限循环；若 $H$ 报告 $M_x$ 在输入~$x$ 上不停机，$G$ 就直接停机。$G$ 以自身的指标作为输入时会停机吗？上述论证表明，它停机当且仅当它不停机——矛盾。因此，存在这样一台Turing机~$H$ 的假设必为假。
\end{explain}
\end{tagblock}

\end{document}